% ============================================================================
% MODEL-SELECTION section — STAGING draft. Model 1 (shared power-budget
% abstraction) locked. Numbers from full 7-model RESISC45 @ n=1300, all Q4.
% Pairs with CAPABILITY_analysis_draft.tex (routing / non-monotonicity).
% ============================================================================

% ---- Calibration + Pareto (model selection) ----
We profile seven vision-language models on RESISC45 (1300 images, 19 to 41 per
category), each at 4-bit quantization on the Jetson Orin NX in its 15\,W power
mode. Fig.~\ref{fig:tradeoff} places every model on the accuracy-energy plane.
Four models lie on the Pareto frontier: Qwen2-VL-2B (0.394 accuracy, 3.93\,J),
Qwen2.5-VL-3B (0.456, 5.72\,J), Qwen2.5-VL-7B (0.549, 9.30\,J), and InternVL3-8B,
which reaches the highest accuracy (0.562) at the highest energy (15.00\,J). The
remaining three are dominated: SmolVLM2-2.2B (0.310), InternVL2.5-4B (0.548), and
gemma-3-4b (0.538, at 37.79\,J).

We adopt the Qwen family as the deployment backbone for an architectural reason,
not a per-point accuracy judgement. Qwen2-VL-2B, Qwen2.5-VL-3B, and
Qwen2.5-VL-7B share one architecture and one activation interface, so they form a
single three-tier ladder that maps directly onto the deployment tiers of this
work: a light single-satellite model, a heavier single-satellite model, and a
split-only model. Their uniform transformer-block structure also keeps the
activation dimension constant at any partition point, which is the precondition
for the two-stage split. InternVL3-8B is an isolated 8B model of a different
architecture: it does not compose with the on-board 2B and 3B tiers into one
routable family, its larger size is even less feasible on a single satellite, and
splitting it is heavier. Its accuracy lead therefore does not change the
deployment choice. InternVL2.5-4B is within statistical noise of Qwen2.5-VL-7B
(0.548 versus 0.549). SmolVLM2-2.2B is the weakest model and is dominated by the
smaller Qwen2-VL-2B, which shows that a model's frontier position is set by
architecture and training rather than by parameter count alone.

% ---- Memory-necessity (why the split) ----
The split pipeline is required by a hardware constraint that does not depend on
any accuracy comparison. A 4-bit Qwen2.5-VL-7B model, together with its vision
encoder and key-value cache, does not fit in the memory left on an 8\,GB
satellite after the on-board sensing pipeline. The two-satellite split is
therefore the only way to run the 7B tier at all, which is what makes the
pipeline necessary rather than optional.

% ---- Peak-power calibration + the shared-budget abstraction ----
The instantaneous power cap is the second structural constraint. We model it as a
shared power budget of $P^{\text{cap}}{=}15$\,W, taking the Jetson 15\,W envelope
as the nominal value; the platform baseline load $P_{\text{base}}$ and the
inference peak draw from this budget. This is an abstraction: a real satellite
separates the platform power system from the compute payload, and a higher-
fidelity power-architecture model is left to future work. We calibrate the peak
bound $\hat{P}_{im}^{\text{CP}}$ as the maximum per-inference peak observed over
the benchmark, a conservative upper bound appropriate for a hard constraint that
must never be violated. Among the deployed configurations, only the 7B tier
shows a sharp transient: its peak reaches 6.37\,W against a 3.77\,W median (a
$1.9\times$ ratio), traced to a sustained compute-bound prefill phase, whereas
the 2B and 3B peaks (3.75 and 4.02\,W) sit close to their own high quantiles. The
7B peak plus a platform baseline in the few-watt range approaches the 15\,W
budget, so a scheduler that commits the split without checking the peak can drive
the combined draw past the cap. This is the empirical basis for the per-slot peak
shield.

% ============================================================================
% Table IV STRUCTURE (numbers pending the max-calibration + nominal-P_base re-run)
% ----------------------------------------------------------------------------
% Rows: Proposed, Greedy-Q, Greedy-E, Random, Static (all draw from the SAME
%   shared feasible set; differ only in selection rule; random tie-break).
% Columns (mean +/- std over 10 seeds):
%   Policy | Peak viol. | Blackout (downtime) | Service rate | 7B splits |
%   Min-fill (max-min) | Total quality
% Reading (to hold after re-run):
%   - Peak viol. = 0 for all (shared feasibility filters infeasible actions);
%     the shield's value shows up in the P_base sweep, not here.
%   - Proposed: 0 blackout, runs the 7B split, competitive min-fill.
%   - Greedy-Q / Random: blackout heavily (undeployable) -> disqualified by SLA.
%   - Greedy-E / Static: safe but never reach the 7B tier.
%   - Report ALL metrics honestly, including where the proposed scheduler
%     trades throughput for sustainability + top-tier reach (NOT a single-metric
%     win; the claim is the deployable + top-tier-capable Pareto corner).
% ============================================================================
